\section{Conclusion}
  \label{sec:conclusion}

  We have proposed a structural interpretation of superconductivity
  as projective phase locking of $w=2$ composite configurations
  within a fiber-based geometric framework.

  The mechanism does not rely on a specific pairing interaction.
  Instead, superconductivity emerges when composite formation
  reduces real-space frustration and global fiber coherence
  becomes energetically favorable.

  In conventional superconductors,
  the projective gap is stabilized by mediator-assisted binding,
  recovering the London and Ginzburg--Landau structure.
  In strongly correlated systems,
  pairing arises from the reduction of staggered frustration,
  naturally separating amplitude and phase scales
  and accommodating pseudogap phenomena.

  The symmetry of the superconducting order parameter
  is determined by a minimization principle
  on the fiber raccordement cost
  under lattice and frustration constraints.
  Applied to nickelates,
  this criterion predicts an extended $s^\pm$ symmetry,
  intermediate disorder sensitivity,
  and a critical temperature scaling
  set by independently measurable exchange and frustration scales.

  The framework is explicitly falsifiable
  through gap symmetry determination,
  phase stiffness scaling,
  disorder response,
  and pressure dependence.

  More broadly,
  the results suggest that superconductivity
  may be understood as a geometric constraint phenomenon:
  when isolated excitations become energetically frustrated,
  composite formation and collective phase locking
  provide the only admissible relaxation channel.

  Further work is required
  to derive the cost functional from microscopic dynamics
  and to extend the analysis to additional material classes.
